\documentclass[10pt,a4paper]{amsart}
\usepackage[margin=19mm]{geometry}
\usepackage{amsmath,amssymb,mathtools}
\usepackage[scr=rsfs,cal=euler]{mathalfa}
\usepackage{xfrac}
\usepackage[unicode,hidelinks,bookmarks=false]{hyperref}
\newcommand{\ie}{i.e.\ }
\newcommand{\cf}{cf.\ }
\newcommand{\resp}{respectively\ }
\newcommand{\Subsection}{\subsection}
\providecommand{\nobreakdash}{\nobreak\hbox{-}}
\setlength{\emergencystretch}{2em}
\multlinegap0pt
\usepackage{bm}
\usepackage{bbm}
\usepackage{dsfont}
\usepackage{xspace}
\usepackage{cancel}
\usepackage{mathtools}
\usepackage{amsthm}
\makeatletter
\@ifundefined{theorem}{\newtheorem{theorem}{Theorem}}{}
\@ifundefined{lemma}{\newtheorem{lemma}[theorem]{Lemma}}{}
\makeatother
\DeclareMathOperator{\den}{den}
\DeclareMathOperator{\num}{num}
\newcommand{\ZZ}{\mathbb{Z}}
\title[Almost all primes are partially regular]{Almost all primes are partially regular}
\author{Evan Chen, Chris Cummins, Ben Eltschig, Dejan Grubisic, Leopold Haller, Letong Hong, Andranik Kurghinyan, Kenny Lau, Hugh Leather, Seewoo Lee, Aram Markosyan, Ken Ono, Manooshree Patel, Gaurang Pendharkar, Vedant Rathi, Alex Schneidman, Volker Seeker, Shubho Sengupta, Ishan Sinha, Jimmy Xin, Jujian Zhang}
\date{Source: arXiv:2602.05090v1 (CC BY 4.0)}
\begin{document}
\maketitle

\noindent\emph{Source and licence.} Almost all primes are partially regular. Version arXiv:2602.05090v1 (CC BY 4.0), distributed under the Creative Commons Attribution 4.0 International licence, as recorded in the arXiv licence metadata for this exact version and shown on the abstract page https://arxiv.org/abs/2602.05090v1. Full rights notice: licenses/arxiv-2602.05090-CC-BY-4.0.txt.\par\smallskip

\noindent\textit{Editorial scope.} Counted unit: the Lemma whose source label is \texttt{lem:num-growth} in the article (source0.tex lines 590--599, source label \texttt{lem:num-growth}), delivered together with the article's own complete original proof of it (source0.tex lines 601--654, one \texttt{proof} environment of 54 lines). No mathematics is written, repaired or re-derived: statement and proof are the article's own text line for line. Required context retained uncounted: lem:euler-zeta (lines 553--559); lem:vsc (lines 565--579). Every definition, display and label of those lines is retained unchanged. Omitted from this package and not redistributed: 1--552, 560--564, 580--589, 600--600, 655--1008. Nothing inside a retained excerpt is omitted or altered: every retained body is delivered exactly as the article's lines are written, so no omission receipt of any kind accompanies this package. Citations: the retained excerpts carry no citation command at all, so no bibliography entry of the article is transcribed and no reference is redistributed. Reference adaptation: none was needed and none was made; every reference inside the delivered unit resolves inside it, so no unresolved reference or citation is produced. Printed slips kept literal: no printed word, symbol, hypothesis or number of the counted statement or of its proof was altered. Layout and class: the article's own class is \texttt{amsart} with packages including amsmath, amssymb, amsthm, bm, mathtools, enumitem, caption, hyperref, graphicx, tikz, comment, geometry, xcolor, amsaddr; the wrapper is the lane's pinned \texttt{amsart} editorial wrapper at 10pt on A4, which restates the theorem-like environments the retained text needs and adds the article's own macro definitions that the retained text needs (\texttt{\textbackslash ZZ}, \texttt{\textbackslash den}, \texttt{\textbackslash num}), with extra wrapper packages (bm, bbm, dsfont, xspace, cancel, mathtools). The delivered numbering is the unit's own. Assets: no retained span contains a figure environment, a graphics-inclusion command, a PDF-page inclusion, a table or a TikZ picture; no image file is copied into this package, the package declares no asset dependency and no image file is redistributed with it. Licence: version arXiv:2602.05090v1 of this article is distributed under the Creative Commons Attribution 4.0 International licence, as recorded in the arXiv licence metadata for this exact version and shown on the abstract page https://arxiv.org/abs/2602.05090v1. A full rights notice is delivered at licenses/arxiv-2602.05090-CC-BY-4.0.txt. Characters: every retained excerpt is pure ASCII, so the delivered engine, XeLaTeX with its default Latin Modern text font, reports no missing character.
\begin{lemma}[Euler's formula for $\zeta(2k)$]\label{lem:euler-zeta}
For each integer $k\ge 1$, we have
\[
\zeta(2k)=(-1)^{k+1}\,\frac{(2\pi)^{2k}}{2(2k)!}\,B_{2k} \ \qquad {\text and}\ \qquad
\zeta(1-2k)=-\frac{B_{2k}}{2k}.
\]
\end{lemma}

\begin{lemma}[von Staudt--Clausen]\label{lem:vsc}
For each integer $k\ge 1$, we have
\[
B_{2k}+\sum_{\substack{\ell\ \mathrm{prime}\\ (\ell-1)\mid 2k}}\frac{1}{\ell}\ \in\ \ZZ.
\]
In particular, we have that
\[
\den(B_{2k})=\prod_{\substack{\ell\ \mathrm{prime}\\ (\ell-1)\mid 2k}}\ell, \label{eqn:lean-mess}
\]
which implies that
\begin{equation}
\label{eqn:vsc_consequence}
    \den(B_{2k})\le \prod_{\ell\le 2k+1}\ell.
\end{equation}
\end{lemma}

\begin{lemma}\label{lem:num-growth}
There is an absolute constant $C_1>0$ such that for every integer $k\ge 1$, we have
\[
\log\bigl|\num(B_{2k})\bigr|\ \le\ C_1\,k\log(2k).
\]
Consequently, there is an absolute constant $C_2>0$ such that for all $m\ge 2$,
\[
\log P_m\ \le\ C_2\,m^2\log m.
\]
\end{lemma}

\begin{proof}
Fix $k\ge 1$.
By Lemma~\ref{lem:euler-zeta}, we have
\[
|B_{2k}|=\frac{2(2k)!\,\zeta(2k)}{(2\pi)^{2k}}
\le \frac{2(2k)!\,\zeta(2)}{(2\pi)^{2k}}
< \frac{4(2k)!}{(2\pi)^{2k}},
\]
since $\zeta(2)=\pi^2/6<2$.
Thanks to Lemma~\ref{lem:vsc}, denominators of Bernoulli numbers are square-free, and satisfy
\[
\den(B_{2k})\le (2k+1)!.
\]
Writing $B_{2k}=\num(B_{2k})/\den(B_{2k})$ in lowest terms gives
$|\num(B_{2k})|=|B_{2k}|\den(B_{2k})$, hence
\[
|\num(B_{2k})|
\le \left(\frac{4(2k)!}{(2\pi)^{2k}}\right) (2k+1)!.
\]
Taking logs and using that $-2k\log(2\pi)<0$, we may discard the negative term to obtain
\[
\log\bigl|\num(B_{2k})\bigr|
\le \log 4 + \log\bigl((2k)!\bigr) + \log\bigl((2k+1)!\bigr).
\]
Using the elementary bound $\log(n!)=\sum_{j=1}^n \log j \le n\log n$, we get
\[
\log\bigl((2k)!\bigr)\le 2k\log(2k)
\quad\text{and}\quad
\log\bigl((2k+1)!\bigr)\le (2k+1)\log(2k+1).
\]
Moreover, since $2k+1\le 4k$, we have $\log(2k+1)\le \log(4k)=\log(2k)+\log 2\le 2\log(2k)$,
and also $2k+1\le 3k$ for $k\ge 1$. Hence
\[
\log\bigl((2k+1)!\bigr)\le (2k+1)\log(2k+1)
\le (2k+1)\log(4k)
\le 3k\cdot 2\log(2k)=6k\log(2k).
\]
Putting these estimates together yields
\[
\log\bigl|\num(B_{2k})\bigr|
\le \log 4 + 2k\log(2k) + 6k\log(2k)
\le C_1\,k\log(2k)
\]
for a suitable absolute constant $C_1>0$.

For the second claim, summing over $1\le k\le m$ and using $\log(2k)\le \log(2m)$ gives
\[
\log P_m=\sum_{k=1}^m \log\bigl|\num(B_{2k})\bigr|
\le C_1 \log(2m)\sum_{k=1}^m k
= \frac{C_1}{2}m(m+1)\log(2m)
\le C_2\,m^2\log m
\]
for all $m\ge 2$ and a suitable absolute constant $C_2>0$.
\end{proof}

\end{document}
